% Generated draft: requires mathematical and source review.
\documentclass[11pt]{article}
\usepackage[margin=25mm]{geometry}
\usepackage{fontspec}
\setmainfont{Latin Modern Roman}
\usepackage{amsmath,amssymb,mathtools,unicode-math}
\setmathfont{Latin Modern Math}
\usepackage{enumitem,needspace,xurl,hyperref}
\setlength{\parindent}{0pt}
\setlength{\parskip}{5pt}
\setlength{\emergencystretch}{3em}
\newcommand{\N}{\mathbb{N}}
\newcommand{\Z}{\mathbb{Z}}
\newcommand{\Q}{\mathbb{Q}}
\newcommand{\R}{\mathbb{R}}
\newcommand{\C}{\mathbb{C}}
\newcommand{\F}{\mathbb{F}}
\newcommand{\A}{\mathbb{A}}
\newcommand{\PP}{\mathbb{P}}
\newcommand{\E}{\mathbb{E}}
\newcommand{\eps}{\varepsilon}
\newcommand{\dd}{\,\mathrm{d}}
\begin{document}
\section*{UnsolvedMath ID: 2301030}
\textbf{Source problem code:} catalogue entry
\subsection{Definitions and mathematical statement}
Let $f$ be a nonconstant meromorphic function on $\C$, with Nevanlinna characteristic $T(r,f)$ and order
\[
\rho(f)=\limsup_{r\to\infty}\frac{\log T(r,f)}{\log r}.
\]
A finite asymptotic value is an $a\in\C$ for which there is a continuous curve $\gamma:[0,\infty)\to\C$ with $|\gamma(t)|\to\infty$ and $f(\gamma(t))\to a$. Let $A_{\mathrm{fin}}(f)$ be the set of distinct such values; values are counted once even if approached through several different tracts.

W. Al-Katifi asks for an upper bound on $\#A_{\mathrm{fin}}(f)$ that takes into account both $\rho(f)$ and the angular measures of the asymptotic tracts. For illustration, a tract over $a$ may be represented by compatible unbounded components of inverse images of small disks $\{w:|w-a|<\eta\}$, and a chosen planar tract $U$ has circular angular size
\[
\theta_U(r)=\operatorname{meas}\{\vartheta\in[0,2\pi):re^{i\vartheta}\in U\}.
\]
The source does not specify this particular tract convention, how shrinking disk components are selected, or which functional of angular size is to enter a bound. The display therefore explains one common meaning of angular measure without prescribing the requested inequality.

A precise bound must specify the admissible angular data and its uniformity across tracts and radii. It must also distinguish all asymptotic tracts from direct or logarithmic tracts: the source does not restrict the singularity type, so such a restriction cannot be inserted silently. Finite order is not explicitly assumed, and no proposed numerical inequality or dependence on $\rho$ and the angular data is supplied. The reported question remains an open-ended quantitative-bound problem. The supplied summary says it was open in Hayman's 2018 edition and that no recent resolution was located.
\subsection{Short English statement}
Al-Katifi asks how to bound the number of distinct finite asymptotic values of a plane meromorphic function using both its growth order and the angular sizes of its tracts. The tract type, angular statistic, and exact form of the desired bound are unspecified.
\subsection{Sources}
\begin{itemize}
\item Local UnsolvedMath registry, permanent ID 2301030, exact source code catalogue entry, ``Research Problems in Function Theory --- Problem 1.30''; status open, classification OPEN-TRIAGE.
\item Hayman, \emph{Research Problems in Function Theory} (2018), Problem 1.30, attributed to W. Al-Katifi, \url{https://arxiv.org/abs/1809.07200}. The record identifies extraction from public arXiv source TeX, accessed July 29, 2026; the 2018 update was retained for status triage.
\item The supplied research summary reports openness as of the 2018 edition and no recent resolution located. It supplies no additional authoritative citation, partial bound, or precise recent search date.
\end{itemize}
\end{document}
